\subsection{Downloading NEMO}
\label{sec:tutor-download}

The official NEMO source code is also managed using the Git revision control system and again there is public read access to its remote repository hosted on GitLab.
To create a local copy, use the following command after changing to the directory in which you wish to download the code:

\begin{verbatim}
  $ git clone https://forge.nemo-ocean.eu/nemo/nemo.git nemo
\end{verbatim}

This will create a subdirectory called \verb|nemo| and download the NEMO source to it.
This \verb|nemo| directory will be referred to as the \textit{top-level source directory} in what follows.

The leading edge of NEMO development takes place on the so-called main branch of its source code repository.
There are also official release versions from time to time.
The source code under the top-level source directory initially corresponds to the latest commit to the main branch.
The instructions in this document have been tested with the 5.0.2 release, although they are likely to work with the main branch as well.
It is recommended that you `checkout' (switch to) this version, which can be done by giving the following command in the top-level source directory:\footnote{
    Git will warn that you are now in `detached HEAD' state: this is because the 5.0.2 release is a specific commit on the branch called \texttt{branch\_5.0}.
    Unless for some reason you are intending (and able) to push changes to the remote repository, this warning can simply be ignored.
}

\begin{verbatim}
  $ git checkout 5.0.2
\end{verbatim}

After giving this command, the source code under the top-level source directory will correspond to the older revision with which the instructions have been tested.
Later, if you wish to return to the head of the main (development) branch, you can do so with the following command:

\begin{verbatim}
  $ git checkout main
\end{verbatim}

Alternatively, you can clone the repository at the point corresponding to release version 5.0.2 straight away (similar to what we did for XIOS in the previous section):

\begingroup\small
\begin{verbatim}
  $ git clone --branch 5.0.2 https://forge.nemo-ocean.eu/nemo/nemo.git nemo
\end{verbatim}
\endgroup
